% ============================================================================
%  styles/contents.tex — 目录版式（书籍形制）
%  原模板是 A4 双栏目录；书籍出版物用单栏、章行带主题色短线的目录。
%
%  实现要点：
%   1. \numberline 用 \@tempdima 作标签槽宽，平时由 \@dottedtocline 设置。
%      自定义 \l@chapter / \l@section 接手后必须自己给 \@tempdima 赋值。
%   2. 不能用 minipage 排版目录行 —— minipage 内部会把 \@tempdima 当作自己的
%      版心宽来用，把上面的槽宽赋值覆盖掉，标签槽变成整行宽，标题整段丢失。
%      这里改用 book 类原生手法：\hangindent + \rightskip + \parfillskip，
%      页码用 \hfill 推到右端，与 \@dottedtocline 同源，稳。
%   3. styles/*.tex 由 lonetrailbook.cls 读入，此时 @ 已是 letter，不需要
%      \makeatletter / \makeatother —— 尾部的 \makeatother 会污染后续读入的文件。
% ============================================================================

\renewcommand{\contentsname}{目录}
\setcounter{tocdepth}{2}   % 章 + 节

% 页码槽宽度（不用 \@pnumwidth，避免与 book 类内部赋值互相覆盖）
\newlength{\LT@pnumw}
\setlength{\LT@pnumw}{2.6em}

% --- 章行：章名 + 页码 + 主题色短线 ----------------------------------------
\renewcommand*\l@chapter[2]{%
  \par\addvspace{12pt}%
  \begingroup
    \setlength{\@tempdima}{2.8em}%   % \numberline 的章号槽宽
    \noindent
    \hangindent\@tempdima
    \hangafter=1
    \rightskip\LT@pnumw
    \parfillskip-\LT@pnumw
    {\bfseries\color{InkBrown}\fontsize{13}{19}\selectfont
      #1\hfill\hb@xt@\LT@pnumw{\hss\bfseries\color{InkBrown}%
        \fontsize{12}{16}\selectfont #2}}\par
  \endgroup
  \nointerlineskip
  \vspace{2.5pt}%
  \noindent{\color{structurecolor}\rule{\linewidth}{0.9pt}}\par
}

% --- 节行：悬挂缩进，安静，无线 --------------------------------------------
\renewcommand*\l@section[2]{%
  \par\addvspace{2.5pt}%
  \begingroup
    \setlength{\@tempdima}{3.4em}%   % \numberline 的节号槽宽
    \noindent
    \hangindent\@tempdima
    \hangafter=1
    \rightskip\LT@pnumw
    \parfillskip-\LT@pnumw
    {\color{InkBrown}\fontsize{10.5}{15}\selectfont
      #1\hfill\hb@xt@\LT@pnumw{\hss\color{InkMuted}%
        \fontsize{10.5}{15}\selectfont #2}}\par
  \endgroup
}
